\documentclass[12pt,a4paper,oneside]{report} % Single-side
%\documentclass[12pt,a4paper,twoside,openright]{report}  % Duplex

% Alapcsomagok betöltése (a sablon része)
\input{include/packages}

% -----------------------------------------------------------------------------
% METAADATOK
% -----------------------------------------------------------------------------
\newcommand{\vikszerzoVezeteknev}{Csabuda}
\newcommand{\vikszerzoKeresztnev}{Nóra}

\newcommand{\vikkonzulensAMegszolitas}{dr.}
\newcommand{\vikkonzulensAVezeteknev}{Martinek}
\newcommand{\vikkonzulensAKeresztnev}{Péter}

\newcommand{\vikcim}{Automatizált számlakezelés SAP Build Process Automation alkalmazásával} % Cím
\newcommand{\viktanszek}{Elektronikai Technológia Tanszék} % Tanszék
\newcommand{\vikdoktipus}{\bsc} % Dokumentum típusa (\bsc vagy \msc)
\newcommand{\vikmunkatipusat}{szakdolgozatot} % a "hallgató nyilatkozat" részhez: szakdolgozatot vagy diplomatervet

\newcommand{\szerzoMeta}{\vikszerzoVezeteknev{} \vikszerzoKeresztnev} % egy szerző esetén

% Beállítások magyar nyelvű dolgozathoz
\input{include/thesis-hu}
\input{include/preamble}

% =============================================================================
% SPECIÁLIS BEÁLLÍTÁSOK (Kódok és Ábrák formázása)
% =============================================================================

% 1. Képek és ábrák rögzítése
\usepackage{float}  % <--- EZ ELENGEDHETETLEN a [H] paraméter használatához!

% 2. Kódblokkok formázása
\usepackage{listings}
\usepackage{xcolor}
\usepackage{pgfplots}

% Színek definiálása a kódkiemeléshez
\definecolor{codegreen}{rgb}{0,0.6,0}
\definecolor{codegray}{rgb}{0.5,0.5,0.5}
\definecolor{codepurple}{rgb}{0.58,0,0.82}
\definecolor{backcolour}{rgb}{0.95,0.95,0.92}

% JavaScript nyelv definíciója (alapból hiányzik a LaTeX-ből)
\lstdefinelanguage{JavaScript}{
  keywords={typeof, new, true, false, catch, function, return, null, catch, switch, var, if, in, while, do, else, case, break},
  keywordstyle=\color{blue}\bfseries,
  ndkeywords={class, export, boolean, throw, implements, import, this},
  ndkeywordstyle=\color{darkgray}\bfseries,
  identifierstyle=\color{black},
  sensitive=false,
  comment=[l]{//},
  morecomment=[s]{/*}{*/},
  commentstyle=\color{codegreen}\ttfamily,
  stringstyle=\color{red}\ttfamily,
  morestring=[b]',
  morestring=[b]"
}

\lstdefinelanguage{Python}{
  keywords={and, as, assert, break, class, continue, def, del, elif, else, except, exec, finally, for, from, global, if, import, in, is, lambda, not, or, pass, print, raise, return, try, while, with, yield, True, False, None},
  keywordstyle=\color{blue}\bfseries,
  ndkeywords={self},
  ndkeywordstyle=\color{darkgray}\bfseries,
  identifierstyle=\color{black},
  sensitive=true,
  comment=[l]{\#},
  morecomment=[s]{"""}{"""},
  commentstyle=\color{codegreen}\ttfamily,
  stringstyle=\color{red}\ttfamily,
  morestring=[b]',
  morestring=[b]"
}

% A listings csomag globális beállításai
\lstset{
    backgroundcolor=\color{backcolour},   
    commentstyle=\color{codegreen},
    keywordstyle=\color{magenta},
    numberstyle=\tiny\color{codegray},
    stringstyle=\color{codepurple},
    basicstyle=\ttfamily\footnotesize, % Kicsit kisebb, írógép betűtípus
    breakatwhitespace=false,         
    breaklines=true,                 
    captionpos=b,                    
    keepspaces=true,                 
    numbers=left,                    
    numbersep=5pt,                  
    showspaces=false,                
    showstringspaces=false,
    showtabs=false,                  
    tabsize=2,
    frame=single                     % Keret a kód köré
}
% =============================================================================


%--------------------------------------------------------------------------------------
% Table of contents and the main text
%--------------------------------------------------------------------------------------
\begin{document}

\pagenumbering{gobble}

\selectthesislanguage

%~~~~~~~~~~~~~~~~~~~~~~~~~~~~~~~~~~~~~~~~~~~~~~~~~~~~~~~~~~~~~~~~~~~~~~~~~~~~~~~~~~~~~~
\include{include/titlepage}        % Szakdolgozat/Diplomaterv címlap


% Table of Contents
%~~~~~~~~~~~~~~~~~~~~~~~~~~~~~~~~~~~~~~~~~~~~~~~~~~~~~~~~~~~~~~~~~~~~~~~~~~~~~~~~~~~~~~
\tableofcontents\cleardoublepage


% Declaration and Abstract
%~~~~~~~~~~~~~~~~~~~~~~~~~~~~~~~~~~~~~~~~~~~~~~~~~~~~~~~~~~~~~~~~~~~~~~~~~~~~~~~~~~~~~~
\include{include/declaration} % Hallgatói nyilatkozat
\pagenumbering{roman}
\setcounter{page}{1}

\selecthungarian

%----------------------------------------------------------------------------
% Abstract in Hungarian
%----------------------------------------------------------------------------
\chapter*{Kivonat}\addcontentsline{toc}{chapter}{Kivonat}

A pénzügyi szektor digitális transzformációjának egyik legnagyobb gátja a mai napig a manuális, papír- vagy PDF-alapú dokumentumfeldolgozás. A vállalati gyakorlatban a számlák jelentős része strukturálatlan formában érkezik, melyek feldolgozása a különböző rendszerek közötti manuális adatátvitellel („swivel-chair integration”) történik. Ez a folyamat nemcsak erőforrás-igényes és lassú, de magas hibakockázatot és átláthatatlanságot is eredményez a beszerzéstől a fizetésig (Purchase-to-Pay) tartó láncban.

Szakdolgozatomban ezen problémára kínálok megoldást a Hiperautomatizálás (Hyperautomation) eszköztárának alkalmazásával. A dolgozat célja egy olyan integrált, „end-to-end” számlafeldolgozó prototípus tervezése és megvalósítása az SAP Business Technology Platform (BTP) felhőalapú környezetében, amely minimalizálja az emberi beavatkozást, ugyanakkor biztosítja a folyamat kontrollálhatóságát.

A megvalósítás során az SAP Build Process Automation (SBPA) technológiát alkalmaztam, ötvözve a robotikus folyamatautomatizálást (RPA) és a mesterséges intelligenciát. A rendszer bemenetét egy saját fejlesztésű, determinisztikus Python szkripttel előállított szintetikus számlaállomány biztosítja, amely lehetővé tette a különböző hibatípusok (hiányos adat, formátumhiba) szisztematikus tesztelését. A strukturálatlan adatok kinyerését az SAP Document Information Extraction (DOX) szolgáltatás legújabb, Generative AI alapú modelljével valósítottam meg. A rendszer architektúrájában kiemelt szerepet kapott az egyedi JavaScript kódokkal vezérelt API kommunikáció és adattranszformáció, amely a technikai robot (Bot) és az üzleti folyamat (Process) közötti hidat képezi. A felhasználói interakciókat és a jóváhagyási döntéseket („Human-in-the-loop”) az SAP Build Work Zone egységesített felületén integráltam.

A prototípus tesztelése három különböző forgatókönyv alapján igazolta, hogy a megoldás képes a számlák teljesen automatikus („touchless”) feldolgozására, valamint a kivételes esetek intelligens kezelésére. Az eredmények azt mutatják, hogy az automatizált folyamat a manuális gyakorlathoz képest radikálisan csökkenti az átfutási időt és az adatrögzítési hibák számát, miközben valós idejű átláthatóságot biztosít a pénzügyi folyamatokban.

\vfill
\selectenglish


%----------------------------------------------------------------------------
% Abstract in English
%----------------------------------------------------------------------------
\chapter*{Abstract}\addcontentsline{toc}{chapter}{Abstract}

One of the biggest obstacles to digital transformation in the financial sector remains manual, paper- or PDF-based document processing. In corporate practice, a significant portion of invoices arrives in unstructured formats, requiring manual data transfer between systems ("swivel-chair integration"). This process is not only resource-intensive and slow but also results in high error risks and a lack of transparency in the Purchase-to-Pay chain.

In my thesis, I propose a solution to this problem by applying the toolkit of Hyperautomation. The aim of the thesis is to design and implement an integrated, end-to-end invoice processing prototype within the cloud-based environment of the SAP Business Technology Platform (BTP), minimizing human intervention while ensuring process control.

For the implementation, I utilized SAP Build Process Automation (SBPA) technology, combining Robotic Process Automation (RPA) with Artificial Intelligence. The system input is provided by a synthetic invoice dataset generated by a self-developed deterministic Python script, which allowed for systematic testing of various error types (missing data, format errors). The extraction of unstructured data was achieved using the latest Generative AI-based model of the SAP Document Information Extraction (DOX) service. A key component of the system architecture is API communication and data transformation driven by custom JavaScript codes, bridging the technical bot and the business process. User interactions and approval decisions ("Human-in-the-loop") were integrated into the unified interface of SAP Build Work Zone.

Testing the prototype across three different scenarios confirmed that the solution is capable of fully automated ("touchless") invoice processing as well as intelligent handling of exception cases. The results demonstrate that the automated process radically reduces turnaround time and data entry errors compared to manual practices, while providing real-time transparency in financial processes.

\vfill
\cleardoublepage

\selectthesislanguage

\newcounter{romanPage}
\setcounter{romanPage}{\value{page}}
\stepcounter{romanPage}    % Összefoglaló


% The main part of the thesis
%~~~~~~~~~~~~~~~~~~~~~~~~~~~~~~~~~~~~~~~~~~~~~~~~~~~~~~~~~~~~~~~~~~~~~~~~~~~~~~~~~~~~~~
\pagenumbering{arabic}

% FEJEZETEK BETÖLTÉSE
\include{content/introduction}
\include{content/theoretical-background}
\include{content/invoice-processing-process-analysis}
\include{content/automated-invoice-processing} % 4. Fejezet (Itt vannak a kódok és ábrák)
\include{content/results-and-evaluation}     % 5. Fejezet
\chapter{Összefoglalás és Következtetések}
\label{chap:summary}

Szakdolgozatomban a pénzügyi folyamatok modernizációjának egyik legkritikusabb területét, a szállítói számlafeldolgozás automatizálását vizsgáltam. A kutatás és a gyakorlati fejlesztés célja annak bizonyítása volt, hogy a hiperautomatizálás eszköztára – különösen az SAP Business Technology Platform (BTP) és az SAP Build Process Automation (BPA) – képes hatékony és vállalati szintű választ adni a manuális adminisztráció okozta üzleti kihívásokra.

\section{A Kutatási Eredmények Összegzése}
A dolgozat elméleti részében feltártam a manuális (,,As-Is'') folyamatok rejtett költségeit és kockázatait, rámutatva a ,,swivel-chair'' integráció és az átláthatatlan jóváhagyási láncok tarthatatlanságára. A gyakorlati megvalósítás során egy működőképes, \textit{End-to-End} prototípust hoztam létre, amely sikeresen integrálta a mesterséges intelligencia alapú adatkinyerést (DOX), a robotikus folyamatautomatizálást (RPA) és az emberi döntéshozatalt támogató munkafolyamatokat (Work Zone).

A tesztelési eredmények (5. fejezet) egyértelműen igazolták a koncepció életképességét: a rendszer képes volt a számlák teljesen automatikus (,,touchless'') feldolgozására, ugyanakkor robusztus módon kezelte a kivételes eseteket és a hibás adatokat is.

\section{A Kutatási Kérdések Megválaszolása}
A dolgozat elején (1.4. fejezet) megfogalmazott kutatási kérdésekre az elvégzett munka alapján az alábbi válaszok adhatók:

\begin{enumerate}
    \item \textbf{Milyen komponensekből áll az SAP BPA, és hogyan kapcsolódnak ezek a hiperautomatizálás trendjeihez?}
    
    A 2. fejezetben bemutattam, hogy az SAP BPA nem egyetlen eszköz, hanem a hiperautomatizálás három pillérének (AI, BPM, RPA) szintézise. A prototípus architektúrája (4.2. fejezet) a gyakorlatban is igazolta ezt az integrációt: a \textit{Document Information Extraction} (AI) biztosította a ,,látást'', a \textit{Process Builder} (BPM) a ,,vezérlést'', míg az \textit{Automation} (RPA) a ,,cselekvést''. Ezen komponensek egységes platformon történő kezelése teszi lehetővé a komplex üzleti folyamatok teljes körű lefedését.

    \item \textbf{Hol keletkeznek a legnagyobb költségek és hibák a manuális folyamatban?}
    
    A 3. fejezet elemzése (,,As-Is'' vs. ,,To-Be'') rámutatott, hogy a legnagyobb veszteségforrás a manuális adatrögzítés és a jóváhagyási folyamat átláthatatlansága. A mérések szerint a manuális adatbevitel számlánként átlagosan 3-5 percet vesz igénybe, és magas a hibakockázata (elütések).

    \item \textbf{Megvalósítható-e SAP BPA-val egy integrált prototípus?}
    
    Igen. A 4. fejezetben részletezett implementáció sikeresen bizonyította, hogy a platform alkalmas a teljes életciklus lefedésére: a számla e-mailben történő beérkezésétől a validáción és jóváhagyáson keresztül egészen az SAP rendszerben történő parkolásig. A megoldás képes volt heterogén (különböző formátumú) bemenetek kezelésére és a \textit{Human-in-the-Loop} interakciók biztosítására.

    \item \textbf{Mennyivel csökkenti a prototípus a feldolgozási időt és a hibákat?}
    
    Az 5. fejezet mérései alapján az automatizált megoldás drasztikus, közel \textbf{80\%-os időmegtakarítást} eredményezett (az átlagos feldolgozási idő 45-60 másodpercre csökkent). Minőségi szempontból a determinisztikus működés és a beépített validációs kapuk teljesen \textbf{kiküszöbölték az adatrögzítési hibákat}, míg a Work Zone bevezetése 100\%-os átláthatóságot biztosított a folyamat státuszáról.
\end{enumerate}

\section{Záró Gondolatok}
A szakdolgozat elkészítése során nemcsak technológiai ismereteket szereztem az SAP legújabb felhőalapú megoldásairól, hanem betekintést nyertem a vállalati folyamatszervezés és a rendszermérnöki tervezés kihívásaiba is. A ,,Citizen Developer'' szemléletmód és a \textit{Low-Code/No-Code} eszközök alkalmazása meggyőzött arról, hogy a jövő mérnökinformatikusának feladata már nem csupán a kódolás, hanem az üzleti értékteremtés és a technológiai lehetőségek kreatív összehangolása.

Meggyőződésem, hogy a dolgozatban bemutatott prototípus nem csupán elméleti kísérlet, hanem egy valós ipari igényre adott, skálázható és továbbfejleszthető mérnöki válasz.

% Opcionális fejezetek (ha nem kellenek, maradhatnak kommentezve)
%\include{content/latex-tools}
%\include{content/thesis-format}
%\include{content/template-usage}


% Acknowledgements
%~~~~~~~~~~~~~~~~~~~~~~~~~~~~~~~~~~~~~~~~~~~~~~~~~~~~~~~~~~~~~~~~~~~~~~~~~~~~~~~~~~~~~~
\include{content/acknowledgement}


% List of Figures, Tables (Opcionális listák)
%~~~~~~~~~~~~~~~~~~~~~~~~~~~~~~~~~~~~~~~~~~~~~~~~~~~~~~~~~~~~~~~~~~~~~~~~~~~~~~~~~~~~~~
%\listoffigures\addcontentsline{toc}{chapter}{\listfigurename}
%\listoftables\addcontentsline{toc}{chapter}{\listtablename}

% Bibliography
%~~~~~~~~~~~~~~~~~~~~~~~~~~~~~~~~~~~~~~~~~~~~~~~~~~~~~~~~~~~~~~~~~~~~~~~~~~~~~~~~~~~~~~
\addcontentsline{toc}{chapter}{\bibname} % Beteszi a tartalomjegyzékbe

\bibliography{bib/mybib} % 3. ADATBÁZIS (ez generálja le a listát)

% Appendix
%~~~~~~~~~~~~~~~~~~~~~~~~~~~~~~~~~~~~~~~~~~~~~~~~~~~~~~~~~~~~~~~~~~~~~~~~~~~~~~~~~~~~~~
%----------------------------------------------------------------------------
\appendix
%----------------------------------------------------------------------------
\chapter*{\fuggelek}\addcontentsline{toc}{chapter}{\fuggelek}
\setcounter{chapter}{\appendixnumber}
%\setcounter{equation}{0} % a fofejezet-szamlalo az angol ABC 6. betuje (F) lesz
\numberwithin{equation}{section}
\numberwithin{figure}{section}
\numberwithin{lstlisting}{section}
%\numberwithin{tabular}{section}

\section{A Tesztadat-generátor Forráskódja (Python)}
\label{app:python_generator}

Az alábbi Python szkript felelős a szintetikus tesztszámlák (PDF) előállításáért. A kód az \texttt{FPDF} könyvtárat használja, és három különböző sablon (\textit{Modern, Simple, General}) alapján képes dinamikusan generálni a dokumentumokat.

\begin{lstlisting}[language=Python, caption={A ComplexInvoice osztály és a PDF generáló logika}, label={code:app_python_full}]
from fpdf import FPDF
import random

class ComplexInvoice(FPDF):
    def header(self):
        # Fejlec generalasa (Logo, Cegadatok)
        self.set_font('Arial', 'B', 12)
        self.cell(0, 10, 'INVOICE', 0, 1, 'C')
        self.ln(10)

    def footer(self):
        # Lablec (Oldalszam)
        self.set_y(-15)
        self.set_font('Arial', 'I', 8)
        self.cell(0, 10, f'Page {self.page_no()}', 0, 0, 'C')

    def generate_invoice(self, data, template_type="General"):
        self.add_page()
        self.set_font('Arial', '', 12)
        
        # Szallito es Vevo adatok
        self.cell(0, 10, f"Sender: {data.get('sender_text')}", 0, 1)
        self.cell(0, 10, f"Bill To: {data.get('bill_to_text')}", 0, 1)
        self.ln(5)
        
        # Szamla reszletei (Tablazatos forma)
        self.set_fill_color(200, 220, 255)
        self.cell(40, 10, 'Description', 1, 0, 'C', 1)
        self.cell(30, 10, 'Quantity', 1, 0, 'C', 1)
        self.cell(30, 10, 'Unit Price', 1, 0, 'C', 1)
        self.cell(30, 10, 'Total', 1, 1, 'C', 1)
        
        for item in data.get('items', []):
            self.cell(40, 10, str(item['desc']), 1)
            self.cell(30, 10, str(item['qty']), 1)
            self.cell(30, 10, str(item['unit']), 1)
            self.cell(30, 10, str(item['total']), 1, 1)
            
        self.ln(10)
        self.set_font('Arial', 'B', 12)
        self.cell(0, 10, f"Total Amount: {data.get('total_gross')} {data.get('currency')}", 0, 1, 'R')

# Pelda hivas:
# pdf = ComplexInvoice()
# pdf.generate_invoice(base_data)
# pdf.output("test_invoice.pdf")
\end{lstlisting}

\section{SAP Build Process Automation Szkriptek}
\label{app:bpa_scripts}

A prototípusban használt egyedi JavaScript kódok, amelyek a folyamatvezérlés (Automation) során futnak le.

\subsection{Adatfeltöltés és API Hívás (Prepare Upload)}

Ez a függvény készíti elő a \texttt{multipart/form-data} kérést a DOX szolgáltatás felé.

\begin{lstlisting}[language=JavaScript, caption={A dokumentumfeltöltést előkészítő JavaScript kód}, label={code:app_js_upload}]
// Input: filePath, accessToken
var myToken = "Bearer " + accessToken;

var optionsPayload = {
    documentType: "custom",
    schemaName: "SAP_invoice_schema",
    clientId: "default"
};

return {
    method: "POST",
    url: "https://aiservices-dox.cfapps.eu10.hana.ondemand.com/document-information-extraction/v1/document/jobs",
    responseType: "json",
    resolveBodyOnly: true,
    metadataType: irpa_core.enums.request.metadataType.formData, 
    headers: {
        Authorization: myToken
    },
    metadata: [
        { name: "file", file: filePath },
        { name: "options", value: JSON.stringify(optionsPayload) }
    ]
};
\end{lstlisting}

\subsection{Adatkinyerés és Normalizálás (Extract Data)}

A DOX szolgáltatásból érkező nyers JSON válasz feldolgozása és egyszerűsített objektummá alakítása.

\begin{lstlisting}[language=JavaScript, caption={A JSON válasz feldolgozása és normalizálása}, label={code:app_js_extract}]
// Input: apiResult
var extracted = {
    VendorName: "Unknown",
    InvoiceNumber: "",
    GrossAmount: 0,
    Currency: "EUR"
};

if (apiResult && apiResult.extraction) {
    var fields = apiResult.extraction.headerFields;
    for (var i = 0; i < fields.length; i++) {
        var name = fields[i].name;
        var value = fields[i].value;
        
        if (name === "senderName") extracted.VendorName = value;
        if (name === "documentNumber") extracted.InvoiceNumber = value;
        if (name === "totalAmount") extracted.GrossAmount = value;
        if (name === "currencyCode") extracted.Currency = value;
    }
}
return extracted;
\end{lstlisting}

\section{Tesztelési Jegyzőkönyv}
\label{app:test_logs}

Az 5. fejezetben bemutatott tesztesetek részletes futási naplója.

\begin{table}[H]
\centering
\begin{tabular}{|l|l|l|l|}
\hline
\textbf{Dátum} & \textbf{Teszteset} & \textbf{Eredmény} & \textbf{Megjegyzés} \\ \hline
2025.10.20 & TC\_01 (Touchless) & SIKERES & Automatikus jovahagyas megtortent (45s) \\ \hline
2025.10.20 & TC\_02 (Missing ID) & SIKERES & Validacios hiba, elutasito urlap generalva \\ \hline
2025.10.21 & TC\_03 (High Value) & SIKERES & Jovahagyasi feladat megjelent a Work Zone-ban \\ \hline
2025.10.21 & TC\_03 (Approval) & SIKERES & Felhasznaloi jovahagyas utan parkolas OK \\ \hline
\end{tabular}
\caption{Összesített tesztelési napló (2025. október)}
\label{tab:app_test_log}
\end{table}

%\label{page:last}
\end{document}